\chapter*{Introduction}
\phantomsection
\addcontentsline{toc}{chapter}{Introduction}

Modern datacenter workloads increasingly combine general-purpose CPUs with accelerators, while also
requiring high-throughput and low-latency access to persistent storage. In many practical systems,
however, accelerator-to-storage communication still relies on host memory as an intermediate staging
area. This approach is functional, but it introduces avoidable data copies when copied from
permanent storage, additional software overhead in case the CPU needs to copy the data itself, and
contention on host resources caused by the increased pressure on the memory bus. As storage
performance and accelerator parallelism continue to grow, these inefficiencies become a significant
system-level bottleneck.

This thesis addresses this problem by proposing and implementing shared CPU-FPGA filesystem access
to an NVMe device using PCI Express peer-to-peer communication. The key idea is to let an FPGA
accelerator exchange data with the NVMe controller directly, while the host CPU participates only
in configuration and standard filesystem operation. To realize this concept, the work combines
hardware and software co-design: a custom FPGA IP core handles NVMe-related PCIe traffic and queue
maintenance on the accelerator side, and a userspace software stack configures devices and exposes a
block-device abstraction in userspace that can be mounted by the operating system.

From the hardware perspective, the thesis begins with \cref{chap:pcie_express}, where the PCIe
standard is introduced as a fundamental basis for the entire work. The chapter also presents the
peer-to-peer communication concept at the core of the proposed implementation and explains how
entities in the host system are configured and addressed. The standardized PCIe interface used in
modern AMD FPGA accelerator cards is then described in \cref{chap:amd_pcie_ip}. The
\emph{Integrated Block for PCIe IP} has specific requirements that must be followed when user logic
processes PCIe traffic.

The NVMe standard that forms the second key element of the implemented solution is described in
\cref{chap:nvme_desc}. The standard has evolved rapidly since its first draft in 2011, and the
fundamentals of its functionality are presented here. The command-queue system and controller
initialization are explained in detail.

\Cref{chap:dma_iuventus} introduces the \emph{DMA Iuventus} module, integrated with the FPGA PCIe
endpoint interfaces. This thesis uses the NDK framework for modular FPGA design and for software
access to the hardware. The DMA module processes request and completion traffic, updates NVMe
doorbells, serves PCIe reads from accelerator-mapped queue and buffer regions, and submits
read/write commands generated for accelerator requests. The design explicitly supports operation
with queue and data structures mapped into accelerator BAR regions, enabling host-bypass movement
between the FPGA and NVMe device.

The software perspective is described in \cref{chap:filesystem}. The thesis uses SPDK with Linux
ublk to create a path from userspace control to a mountable block device. This allows the solution
to demonstrate shared storage access in a standard Linux environment rather than only raw
sector-level accesses. The selected \emph{exFAT} filesystem provides a simple and portable target for
validating the correctness of read/write operations across CPU and accelerator IO.

The implemented solution is evaluated in \cref{chap:eval} on an \emph{AMD Alveo U55C} accelerator
and an NVMe SSD in a custom server platform. Evaluation includes FPGA resource utilization and
request-latency measurements for multiple access patterns and transfer sizes. The measurements
confirm latency behavior consistent with typical NVMe device operation. The last section of this
chapter describes how the provided application is run and how to access the shared filesystem.

Overall, this thesis demonstrates that shared CPU-FPGA filesystem access over NVMe can be achieved
without mandatory host-memory data staging by combining PCIe peer-to-peer transport, userspace
device control, and an FPGA-side command-processing engine.